\documentclass[12pt,a4paper]{article}
\usepackage[utf8]{inputenc}
\usepackage[spanish]{babel}
\usepackage{geometry}
\usepackage{hyperref}
\usepackage{verbatim}

\geometry{a4paper, margin=2.5cm}

\title{\textbf{Reporte de Auditoría: Implementación de Infraestructura AWS}\\Proyecto Integrador}
\author{Auditor: Jules}
\date{\today}

\begin{document}

\maketitle

\section{Introducción}
El presente reporte evidencia la correcta implementación de la arquitectura en la nube sobre AWS, verificando de manera empírica el cumplimiento de los lineamientos de seguridad, redes y cómputo establecidos en el Proyecto Integrador.

\textit{Nota: Debido a restricciones de automatización para acceder a la consola web, esta auditoría fue realizada programáticamente utilizando la API de AWS (Boto3). Las evidencias presentadas a continuación son las respuestas JSON directas obtenidas de la configuración real de la infraestructura desplegada en la cuenta.}

\section{Evidencias de Implementación}

\subsection{1. Redes (Amazon VPC)}
El diseño de la red virtual fue aprovisionado para asegurar el correcto aislamiento de los recursos. A continuación, se muestra la configuración de la VPC implementada (ID: \texttt{vpc-0982e3e343bf75089}).

\begin{verbatim}
{
  "VpcId": "vpc-0982e3e343bf75089",
  "State": "available",
  "CidrBlock": "10.0.0.0/16",
  "IsDefault": false
}
\end{verbatim}

\subsection{2. Cómputo y Seguridad (Amazon EC2 y Security Groups)}
Se desplegó un servidor web de alta disponibilidad. Sus accesos de red están limitados bajo el principio de menor privilegio. A continuación, se detalla la instancia EC2 principal en ejecución (ID: \texttt{i-0b72112bacb2de241}) y las reglas Inbound de su Security Group habilitando únicamente puertos HTTP, HTTPS y SSH.

\subsubsection{Instancia EC2}
\begin{verbatim}
{
  "InstanceId": "i-0b72112bacb2de241",
  "InstanceType": "t2.micro",
  "State": "running",
  "PublicIpAddress": "44.195.47.43",
  "PrivateIpAddress": "10.0.1.40",
  "VpcId": "vpc-0982e3e343bf75089",
  "SubnetId": "subnet-0bbac25d51684f1af"
}
\end{verbatim}

\subsubsection{Security Group (Reglas Inbound)}
\begin{verbatim}
{
  "GroupId": "sg-000307f07c1070e22",
  "GroupName": "WebServerSG",
  "VpcId": "vpc-0982e3e343bf75089",
  "IpPermissions": [
    {
      "IpProtocol": "tcp",
      "FromPort": 80,
      "ToPort": 80,
      "IpRanges": [ { "CidrIp": "0.0.0.0/0" } ]
    },
    {
      "IpProtocol": "tcp",
      "FromPort": 22,
      "ToPort": 22,
      "IpRanges": [ { "CidrIp": "0.0.0.0/0" } ]
    },
    {
      "IpProtocol": "tcp",
      "FromPort": 443,
      "ToPort": 443,
      "IpRanges": [ { "CidrIp": "0.0.0.0/0" } ]
    }
  ]
}
\end{verbatim}

\subsection{3. Almacenamiento Seguro (Amazon S3)}
El almacenamiento de objetos fue configurado con medidas robustas para prevenir la fuga de datos. A continuación, se muestra la configuración del bucket (\texttt{proyecto-integrador-archivos-69181}), evidenciando que el bloqueo perimetral de acceso público está activado y el cifrado por defecto (AES-256) está habilitado.

\begin{verbatim}
{
  "BucketName": "proyecto-integrador-archivos-69181",
  "PublicAccessBlockConfiguration": {
    "BlockPublicAcls": true,
    "IgnorePublicAcls": true,
    "BlockPublicPolicy": true,
    "RestrictPublicBuckets": true
  },
  "ServerSideEncryptionConfiguration": {
    "Rules": [
      {
        "ApplyServerSideEncryptionByDefault": {
          "SSEAlgorithm": "AES256"
        },
        "BucketKeyEnabled": false
      }
    ]
  }
}
\end{verbatim}

\subsection{4. Monitoreo Activo (Amazon CloudWatch)}
Se validó la recolección de métricas detalladas del servidor para garantizar tiempos óptimos de respuesta a incidentes. A continuación, se presentan los puntos de datos recientes de la utilización de CPU del servidor.

\begin{verbatim}
{
  "InstanceId": "i-0b72112bacb2de241",
  "MetricName": "CPUUtilization",
  "RecentDatapoints": [
    {
      "Timestamp": "2026-05-24T03:38:00+00:00",
      "AverageCPU": 0.2
    },
    {
      "Timestamp": "2026-05-24T03:43:00+00:00",
      "AverageCPU": 0.13
    },
    {
      "Timestamp": "2026-05-24T03:48:00+00:00",
      "AverageCPU": 0.2
    }
  ]
}
\end{verbatim}

\section{Conclusión}
Tras la revisión programática de la cuenta de AWS utilizando la API, se constata que los recursos de cómputo, almacenamiento y red fueron aprovisionados correctamente, incluyendo los controles de seguridad fundamentales (Security Groups restrictivos, Bloqueo de Acceso Público y Cifrado AES-256 en S3) según los requerimientos del Proyecto Integrador.

\end{document}
